%%%%%
%
\section{Humidit� relative � la glace � partir de T et Td}

\itemstyle{Op\'erateur~:} \textsf{HumI\_T\_Td}
\index{HumI\_T\_Td}

%%%%%OBJET%%%%
\itemstyle{Objet~:}
Calcule l'humidit� relative \`a partir de la temp\'erature 
du point de ros\'ee et de la temp\'erature de l'air, en utilisant les
formules relatives \`a la glace. 

%%%%%ARGUMENTS%%%%
\bigskip
%%%%%ENTREES%%%%
\begin{tabarg}
\multicolumn{3}{l}{\titretab{ Entr\'ees}}\\
\hline
$T_s$ & Vecteur, \texttt{float} & Temp\'erature statique de l'air ($K$ ou $^\circ C$) \\
$T_d$ & Vecteur, \texttt{float} & Temp\'erature du point de ros\'ee ($K$ ou $^\circ C$) 
\end{tabarg}

\medskip
%%%%%SORTIES%%%%
\begin{tabarg}
\multicolumn{3}{l}{\titretab{Sorties}}\\
\hline
$H_u$ & Vecteur, \texttt{float} & Humidit\'e relative � la glace ($\%$)\\
\end{tabarg}


%%%%%FORMULE%%%%
\itemstyle{Formule}

%%%% THERMODYNAMIQUE
\[
H_u = 100 \times \frac{e_w(T_d)}{e_{wi}(T_s)} 
\]
Si $H_u > 115 \%$, $Hu = 115 \%$.



%%%%COMMENTAIRE%%%%
\itemstyle{Commentaire~:}

Les fonctions $ew()$ et  $ewi()$ sont d\'ecrites en annexe
\ref{annexe-thermo}.

Pour les mesures m�t�orologiques op�rationnelles (par exemple
les radiosondages), l'humidit� relative est toujours d�finie
par rapport � l'eau liquide, m�me � des temp�ratures
inf�rieures � 0�C (cf. le \emph{Guide des instruments et des m�thodes
d'observation m�t�orologique} \cite{OMM-N8-1996}). Le r�sultat
de cet op�rateur ne peut donc pas �tre directement 
compar� � de telles mesures.

%%%%%AUTEUR%%%%
\itemstyle{Auteur~:}
CNRM/GMEI/TRAMM

%%%%%REFERENCE%%%%
\itemstyle{R\'ef\'erences~:}
